% Chapter 20 infographic -- use right after \label{ch:20}: % Chapter 19 infographic -- use right after \label{ch:19}: % Chapter 18 infographic -- use right after \label{ch:18}: % Chapter 17 infographic -- use right after \label{ch:17}: \input{chapters/ch17-infographic}
\ifcsname FEinfoLoaded\endcsname\else\input{common/fe-infographic}\fi
\begin{FEinfographic}{ch:17}{a trader who sells an option and delta-hedges it as Black--Scholes prescribes is left with little risk}

\infopanel{1}{The Greeks behind the hedge}{%
\infocards
  {Delta}{How much the option's value moves with the stock; holding delta shares offsets small moves.}
  {Gamma}{How fast delta changes, so the hedge must be adjusted and loses on large moves.}
  {Theta}{The value lost as time passes, the seller's pay for bearing gamma.}}

\infopanel{2}{Evidence scoreboard}{%
\begin{infoscore}
\scoreitem{y}{Yes}{Hedging when the model holds}
\scoreitem{n}{No}{Hedging with trading costs}
\scoreitem{n}{No}{Market prices consistent with the model}
\scoreitem{n}{No}{Hedged option positions earn zero}
\scoreitem{p}{Partly}{Better hedges}
\scoreitem{n}{No}{Dynamic replication in a crash}
\end{infoscore}}

\infopanel{3}{Key numbers}{%
\infonumbers
  {about 1/5}{of the premium left as error with daily hedging}
  {4$\times$}{as frequent hedging halves the error}
  {\$90 bn}{at most, covered by portfolio insurance in 1987}
  {over 20\%}{market fall on October 19, 1987}}

\infopanel{4}{Lab: the Black--Scholes hedge along real and simulated paths}{%
\begin{infolab}
\item \textbf{A.} The hedged P\&L of a one-month call at several hedging frequencies, on real and simulated paths.
\item \textbf{B.} The P\&L against the gamma--theta approximation and the implied-minus-realized variance gap.
\item \textbf{C.} Hedging frequency and costs. \textbf{D.} The worst months and their share of all losses.
\end{infolab}}

\infoverdict
  {Delta hedging removes most of the everyday risk of an option position, but what remains is a bet on realized against implied volatility, with large losses in jumps and volatility spikes that more frequent hedging cannot remove.}
  {A fund that sells options and delta-hedges is not market-neutral in the sense that matters: it bets against volatility and loses heavily in jumps.}
\end{FEinfographic}

\ifcsname FEinfoLoaded\endcsname\else% ======================================================================
% fe-infographic.tex -- shared layout for the one-page chapter infographics
% Loaded automatically by each chapters/chNN-infographic.tex.
% Works with the Glory style and the agentic style; colours are provided
% here if the preamble does not define them.
% ======================================================================
\providecolor{FEMaroon}{RGB}{128,30,45}
\providecolor{FECream}{RGB}{253,248,238}
\providecolor{FEGold}{RGB}{196,150,60}
\providecolor{fenavy}{HTML}{1B2A6B}
\providecolor{fegray}{HTML}{4D4D4D}
\providecolor{igRisk}{HTML}{1F4E79}
\providecolor{igLuck}{HTML}{8A8A8A}
\providecolor{igSkill}{HTML}{2E6B3F}
\providecolor{igYes}{HTML}{2E6B3F}
\providecolor{igNo}{HTML}{A33A3A}
\providecolor{igPart}{HTML}{C47F00}
\providecolor{igOpen}{HTML}{5D6D7E}

\def\FEinfoLoaded{1}

% Title band and page frame -----------------------------------------
\newenvironment{FEinfographic}[2]{% #1 = chapter label, #2 = claim
  \clearpage\thispagestyle{plain}%
  \begingroup\sffamily\footnotesize\setlength{\parskip}{0pt}\setlength{\parindent}{0pt}%
  \linespread{1.0}\selectfont
  \begin{tcolorbox}[enhanced, colback=FEMaroon, colframe=FEMaroon, arc=3pt,
      boxrule=0pt, left=8pt, right=8pt, top=5pt, bottom=5pt, before skip=0pt, after skip=5pt]
    {\color{white}\bfseries\large Chapter~\ref{#1} at a glance\par}
    \vspace{2pt}{\color{white}\small Claim on trial: \emph{#2}\par}
  \end{tcolorbox}}
  {\par\endgroup\clearpage}

% A numbered panel ----------------------------------------------------
\newcommand{\infopanel}[3]{% #1 number, #2 heading, #3 body
  \begin{tcolorbox}[enhanced, colback=FECream, colframe=FEGold, boxrule=0.6pt, arc=3pt,
      left=7pt, right=7pt, top=4pt, bottom=5pt, before skip=0pt, after skip=5pt]
    {\color{FEMaroon}\bfseries\small #1\quad #2\par}\vspace{3pt}%
    #3
  \end{tcolorbox}}

% Three idea cards ------------------------------------------------------
\newcommand{\infocard}[3]{% #1 colour, #2 title, #3 text
  \begin{minipage}[t]{0.315\linewidth}
    {\setlength{\fboxsep}{2.5pt}\colorbox{#1}{\makebox[\dimexpr\linewidth-2\fboxsep]{\color{white}\bfseries\rule[-0.3ex]{0pt}{2.2ex}#2}}}\par\vspace{3pt}
    {\scriptsize #3\par}
  \end{minipage}}
\newcommand{\infocards}[6]{%
  \infocard{igRisk}{#1}{#2}\hfill\infocard{igLuck}{#3}{#4}\hfill\infocard{igSkill}{#5}{#6}\par}

% Evidence scoreboard ---------------------------------------------------
\newcommand{\vchip}[2]{% #1 = y/n/p/o, #2 = label
  \edef\FEchipc{\ifx y#1igYes\else\ifx n#1igNo\else\ifx p#1igPart\else igOpen\fi\fi\fi}%
  \tcbox[on line, arc=2pt, boxrule=0pt, left=2pt, right=2pt, top=1pt, bottom=1pt,
    colback=\FEchipc,
    fontupper=\sffamily\bfseries\tiny\color{white}]{\makebox[3.6em]{#2}}}
% Items are laid out two to a row.
\newcount\FEscorecount
\newenvironment{infoscore}{\global\FEscorecount=0 \noindent\ignorespaces}{\par}
\newcommand{\scoreitem}[3]{% #1 class, #2 chip label, #3 situation
  \global\advance\FEscorecount by 1
  \begin{minipage}[t]{0.485\linewidth}\raggedright
    \vchip{#1}{#2}\hspace{4pt}\parbox[t]{\dimexpr\linewidth-5.4em}{\scriptsize #3}\vspace{3pt}
  \end{minipage}%
  \ifodd\FEscorecount\hfill\else\par\noindent\fi\ignorespaces}

% Four key numbers ------------------------------------------------------
\newcount\FEnumgroup
\newcommand{\infonum}[2]{%
  \begin{minipage}[t]{0.235\linewidth}
  \begin{tcolorbox}[enhanced, colback=white, colframe=FEGold, boxrule=0.5pt, arc=2pt,
      left=3pt, right=3pt, top=3pt, bottom=3pt, before skip=0pt, after skip=0pt,
      equal height group=FEnum\the\FEnumgroup, valign=top, halign=center]
    {\color{FEMaroon}\bfseries\Large #1\par}\vspace{2pt}{\scriptsize #2\par}
  \end{tcolorbox}
  \end{minipage}}
\newcommand{\infonumbers}[8]{%
  \global\advance\FEnumgroup by 1
  \infonum{#1}{#2}\hfill\infonum{#3}{#4}\hfill\infonum{#5}{#6}\hfill\infonum{#7}{#8}\par}

% Lab parts -------------------------------------------------------------
\newenvironment{infolab}{\scriptsize\begin{itemize}[leftmargin=1.1em, itemsep=1pt, topsep=1pt, parsep=0pt]}{\end{itemize}}

% Verdict band ----------------------------------------------------------
\newcommand{\infoverdict}[2]{% #1 = where the evidence stands, #2 = what investors can use
  \begin{tcolorbox}[enhanced, colback=white, colframe=FEMaroon, boxrule=0.8pt, arc=3pt,
      leftrule=4pt, left=7pt, right=7pt, top=4pt, bottom=5pt, before skip=0pt, after skip=0pt]
    {\color{FEMaroon}\bfseries\small Where the evidence stands\par}\vspace{2pt}
    {\scriptsize #1\par}\vspace{4pt}
    {\color{FEMaroon}\bfseries\small What investors can use\par}\vspace{2pt}
    {\scriptsize #2\par}
  \end{tcolorbox}}
\fi
\begin{FEinfographic}{ch:18}{is a trader who sells an option and delta-hedges it as Black--Scholes prescribes left with little risk?}

\infopanel{1}{The Greeks behind the hedge}{%
\infocards
  {Delta}{How much the option's value moves with the stock; holding delta shares offsets small moves.}
  {Gamma}{How fast delta changes, so the hedge must be adjusted and loses on large moves.}
  {Theta}{The value lost as time passes, the seller's pay for bearing gamma.}}

\infopanel{2}{Where the evidence stands: what published studies found}{%
\begin{infoscore}
\scoreitem{y}{Yes}{Hedging when the model holds}
\scoreitem{n}{No}{Hedging with trading costs}
\scoreitem{n}{No}{Market prices consistent with the model}
\scoreitem{n}{No}{Hedged option positions earn zero}
\scoreitem{p}{Partly}{Better hedges}
\scoreitem{n}{No}{Dynamic replication in a crash}
\end{infoscore}}

\infopanel{3}{Key numbers}{%
\infonumbers
  {about 1/5}{of the premium left as error with daily hedging}
  {4$\times$}{as frequent hedging halves the error}
  {\$90 bn}{at most, covered by portfolio insurance in 1987}
  {over 20\%}{market fall on October 19, 1987}}

\infopanel{4}{Lab: Delta hedging along real and simulated paths}{%
\begin{infolab}
\item \textbf{A.} The hedged P\&L of a one-month call at several hedging frequencies, on real and simulated paths.
\item \textbf{B.} The P\&L against the gamma--theta approximation and the implied-minus-realized variance gap.
\item \textbf{C.} Hedging frequency and costs. \textbf{D.} The worst months and their share of all losses.
\end{infolab}}

\infosummary
  {Delta hedging removes most of the everyday risk of an option position, but what remains is a bet on realized against implied volatility, with large losses in jumps and volatility spikes that more frequent hedging cannot remove.}
  {A fund that sells options and delta-hedges is not market-neutral in the sense that matters: it bets against volatility and loses heavily in jumps.}
\end{FEinfographic}

\ifcsname FEinfoLoaded\endcsname\else% ======================================================================
% fe-infographic.tex -- shared layout for the one-page chapter infographics
% Loaded automatically by each chapters/chNN-infographic.tex.
% Works with the Glory style and the agentic style; colours are provided
% here if the preamble does not define them.
% ======================================================================
\providecolor{FEMaroon}{RGB}{128,30,45}
\providecolor{FECream}{RGB}{253,248,238}
\providecolor{FEGold}{RGB}{196,150,60}
\providecolor{fenavy}{HTML}{1B2A6B}
\providecolor{fegray}{HTML}{4D4D4D}
\providecolor{igRisk}{HTML}{1F4E79}
\providecolor{igLuck}{HTML}{8A8A8A}
\providecolor{igSkill}{HTML}{2E6B3F}
\providecolor{igYes}{HTML}{2E6B3F}
\providecolor{igNo}{HTML}{A33A3A}
\providecolor{igPart}{HTML}{C47F00}
\providecolor{igOpen}{HTML}{5D6D7E}

\def\FEinfoLoaded{1}

% Title band and page frame -----------------------------------------
\newenvironment{FEinfographic}[2]{% #1 = chapter label, #2 = claim
  \clearpage\thispagestyle{plain}%
  \begingroup\sffamily\footnotesize\setlength{\parskip}{0pt}\setlength{\parindent}{0pt}%
  \linespread{1.0}\selectfont
  \begin{tcolorbox}[enhanced, colback=FEMaroon, colframe=FEMaroon, arc=3pt,
      boxrule=0pt, left=8pt, right=8pt, top=5pt, bottom=5pt, before skip=0pt, after skip=5pt]
    {\color{white}\bfseries\large Chapter~\ref{#1} at a glance\par}
    \vspace{2pt}{\color{white}\small Claim on trial: \emph{#2}\par}
  \end{tcolorbox}}
  {\par\endgroup\clearpage}

% A numbered panel ----------------------------------------------------
\newcommand{\infopanel}[3]{% #1 number, #2 heading, #3 body
  \begin{tcolorbox}[enhanced, colback=FECream, colframe=FEGold, boxrule=0.6pt, arc=3pt,
      left=7pt, right=7pt, top=4pt, bottom=5pt, before skip=0pt, after skip=5pt]
    {\color{FEMaroon}\bfseries\small #1\quad #2\par}\vspace{3pt}%
    #3
  \end{tcolorbox}}

% Three idea cards ------------------------------------------------------
\newcommand{\infocard}[3]{% #1 colour, #2 title, #3 text
  \begin{minipage}[t]{0.315\linewidth}
    {\setlength{\fboxsep}{2.5pt}\colorbox{#1}{\makebox[\dimexpr\linewidth-2\fboxsep]{\color{white}\bfseries\rule[-0.3ex]{0pt}{2.2ex}#2}}}\par\vspace{3pt}
    {\scriptsize #3\par}
  \end{minipage}}
\newcommand{\infocards}[6]{%
  \infocard{igRisk}{#1}{#2}\hfill\infocard{igLuck}{#3}{#4}\hfill\infocard{igSkill}{#5}{#6}\par}

% Evidence scoreboard ---------------------------------------------------
\newcommand{\vchip}[2]{% #1 = y/n/p/o, #2 = label
  \edef\FEchipc{\ifx y#1igYes\else\ifx n#1igNo\else\ifx p#1igPart\else igOpen\fi\fi\fi}%
  \tcbox[on line, arc=2pt, boxrule=0pt, left=2pt, right=2pt, top=1pt, bottom=1pt,
    colback=\FEchipc,
    fontupper=\sffamily\bfseries\tiny\color{white}]{\makebox[3.6em]{#2}}}
% Items are laid out two to a row.
\newcount\FEscorecount
\newenvironment{infoscore}{\global\FEscorecount=0 \noindent\ignorespaces}{\par}
\newcommand{\scoreitem}[3]{% #1 class, #2 chip label, #3 situation
  \global\advance\FEscorecount by 1
  \begin{minipage}[t]{0.485\linewidth}\raggedright
    \vchip{#1}{#2}\hspace{4pt}\parbox[t]{\dimexpr\linewidth-5.4em}{\scriptsize #3}\vspace{3pt}
  \end{minipage}%
  \ifodd\FEscorecount\hfill\else\par\noindent\fi\ignorespaces}

% Four key numbers ------------------------------------------------------
\newcount\FEnumgroup
\newcommand{\infonum}[2]{%
  \begin{minipage}[t]{0.235\linewidth}
  \begin{tcolorbox}[enhanced, colback=white, colframe=FEGold, boxrule=0.5pt, arc=2pt,
      left=3pt, right=3pt, top=3pt, bottom=3pt, before skip=0pt, after skip=0pt,
      equal height group=FEnum\the\FEnumgroup, valign=top, halign=center]
    {\color{FEMaroon}\bfseries\Large #1\par}\vspace{2pt}{\scriptsize #2\par}
  \end{tcolorbox}
  \end{minipage}}
\newcommand{\infonumbers}[8]{%
  \global\advance\FEnumgroup by 1
  \infonum{#1}{#2}\hfill\infonum{#3}{#4}\hfill\infonum{#5}{#6}\hfill\infonum{#7}{#8}\par}

% Lab parts -------------------------------------------------------------
\newenvironment{infolab}{\scriptsize\begin{itemize}[leftmargin=1.1em, itemsep=1pt, topsep=1pt, parsep=0pt]}{\end{itemize}}

% Verdict band ----------------------------------------------------------
\newcommand{\infoverdict}[2]{% #1 = where the evidence stands, #2 = what investors can use
  \begin{tcolorbox}[enhanced, colback=white, colframe=FEMaroon, boxrule=0.8pt, arc=3pt,
      leftrule=4pt, left=7pt, right=7pt, top=4pt, bottom=5pt, before skip=0pt, after skip=0pt]
    {\color{FEMaroon}\bfseries\small Where the evidence stands\par}\vspace{2pt}
    {\scriptsize #1\par}\vspace{4pt}
    {\color{FEMaroon}\bfseries\small What investors can use\par}\vspace{2pt}
    {\scriptsize #2\par}
  \end{tcolorbox}}
\fi
\begin{FEinfographic}{ch:19}{does selling index options earn more than the risk taken, because implied volatility exceeds the volatility that follows?}

\infopanel{1}{Two ways to sell options}{%
\infocards
  {Put-writing}{Sell index puts backed by Treasury bills, collecting a premium in return for crash risk.}
  {Covered calls}{Hold the index and sell calls, giving up gains above the strike for income.}
  {The premium}{Implied volatility has exceeded realized volatility on average: the variance risk premium.}}

\infopanel{2}{Where the evidence stands: what published studies found}{%
\begin{infoscore}
\scoreitem{y}{Yes}{Implied above realized variance}
\scoreitem{y}{Yes}{Option buyers pay for volatility}
\scoreitem{y}{Yes}{An economic reason for the premium}
\scoreitem{p}{Partly}{Covered calls against the index}
\scoreitem{p}{Weakly}{Covered calls after adjusting for beta}
\scoreitem{p}{Weakly}{Risk adjustment that prices skewness}
\scoreitem{n}{No}{Mispricing, judged from put returns}
\end{infoscore}}

\infopanel{3}{Key numbers}{%
\infonumbers
  {about 3\%}{a week lost by zero-beta straddles}
  {1--3\%}{one month's premium, the only cushion in a crash}
  {0.5--0.7}{typical beta of an option-income fund}
  {over 61\%}{of S\&P 500 option volume in same-day options, May 2025}}

\infopanel{4}{Lab: The variance risk premium and option-selling strategies}{%
\begin{infolab}
\item \textbf{A.} The VIX against the realized volatility that followed, month by month.
\item \textbf{B.} An at-the-money put-write and a covered call, against the market.
\item \textbf{C.} Crises, up and down betas, and the MPPM. \textbf{D.} Beta exposure against the pure volatility bet.
\end{infolab}}

\infosummary
  {Implied volatility has been higher than realized volatility on average, so the variance risk premium is real, but covered-call and put-writing strategies earn most of their return from stock-market exposure and lose most of what the market loses in crashes.}
  {Payouts are not yield: compare an option-income fund with an index-and-bills mix of the same beta, and look at its worst months, not its Sharpe ratio.}
\end{FEinfographic}

\ifcsname FEinfoLoaded\endcsname\else% ======================================================================
% fe-infographic.tex -- shared layout for the one-page chapter infographics
% Loaded automatically by each chapters/chNN-infographic.tex.
% Works with the Glory style and the agentic style; colours are provided
% here if the preamble does not define them.
% ======================================================================
\providecolor{FEMaroon}{RGB}{128,30,45}
\providecolor{FECream}{RGB}{253,248,238}
\providecolor{FEGold}{RGB}{196,150,60}
\providecolor{fenavy}{HTML}{1B2A6B}
\providecolor{fegray}{HTML}{4D4D4D}
\providecolor{igRisk}{HTML}{1F4E79}
\providecolor{igLuck}{HTML}{8A8A8A}
\providecolor{igSkill}{HTML}{2E6B3F}
\providecolor{igYes}{HTML}{2E6B3F}
\providecolor{igNo}{HTML}{A33A3A}
\providecolor{igPart}{HTML}{C47F00}
\providecolor{igOpen}{HTML}{5D6D7E}

\def\FEinfoLoaded{1}

% Title band and page frame -----------------------------------------
\newenvironment{FEinfographic}[2]{% #1 = chapter label, #2 = claim
  \clearpage\thispagestyle{plain}%
  \begingroup\sffamily\footnotesize\setlength{\parskip}{0pt}\setlength{\parindent}{0pt}%
  \linespread{1.0}\selectfont
  \begin{tcolorbox}[enhanced, colback=FEMaroon, colframe=FEMaroon, arc=3pt,
      boxrule=0pt, left=8pt, right=8pt, top=5pt, bottom=5pt, before skip=0pt, after skip=5pt]
    {\color{white}\bfseries\large Chapter~\ref{#1} at a glance\par}
    \vspace{2pt}{\color{white}\small Claim on trial: \emph{#2}\par}
  \end{tcolorbox}}
  {\par\endgroup\clearpage}

% A numbered panel ----------------------------------------------------
\newcommand{\infopanel}[3]{% #1 number, #2 heading, #3 body
  \begin{tcolorbox}[enhanced, colback=FECream, colframe=FEGold, boxrule=0.6pt, arc=3pt,
      left=7pt, right=7pt, top=4pt, bottom=5pt, before skip=0pt, after skip=5pt]
    {\color{FEMaroon}\bfseries\small #1\quad #2\par}\vspace{3pt}%
    #3
  \end{tcolorbox}}

% Three idea cards ------------------------------------------------------
\newcommand{\infocard}[3]{% #1 colour, #2 title, #3 text
  \begin{minipage}[t]{0.315\linewidth}
    {\setlength{\fboxsep}{2.5pt}\colorbox{#1}{\makebox[\dimexpr\linewidth-2\fboxsep]{\color{white}\bfseries\rule[-0.3ex]{0pt}{2.2ex}#2}}}\par\vspace{3pt}
    {\scriptsize #3\par}
  \end{minipage}}
\newcommand{\infocards}[6]{%
  \infocard{igRisk}{#1}{#2}\hfill\infocard{igLuck}{#3}{#4}\hfill\infocard{igSkill}{#5}{#6}\par}

% Evidence scoreboard ---------------------------------------------------
\newcommand{\vchip}[2]{% #1 = y/n/p/o, #2 = label
  \edef\FEchipc{\ifx y#1igYes\else\ifx n#1igNo\else\ifx p#1igPart\else igOpen\fi\fi\fi}%
  \tcbox[on line, arc=2pt, boxrule=0pt, left=2pt, right=2pt, top=1pt, bottom=1pt,
    colback=\FEchipc,
    fontupper=\sffamily\bfseries\tiny\color{white}]{\makebox[3.6em]{#2}}}
% Items are laid out two to a row.
\newcount\FEscorecount
\newenvironment{infoscore}{\global\FEscorecount=0 \noindent\ignorespaces}{\par}
\newcommand{\scoreitem}[3]{% #1 class, #2 chip label, #3 situation
  \global\advance\FEscorecount by 1
  \begin{minipage}[t]{0.485\linewidth}\raggedright
    \vchip{#1}{#2}\hspace{4pt}\parbox[t]{\dimexpr\linewidth-5.4em}{\scriptsize #3}\vspace{3pt}
  \end{minipage}%
  \ifodd\FEscorecount\hfill\else\par\noindent\fi\ignorespaces}

% Four key numbers ------------------------------------------------------
\newcount\FEnumgroup
\newcommand{\infonum}[2]{%
  \begin{minipage}[t]{0.235\linewidth}
  \begin{tcolorbox}[enhanced, colback=white, colframe=FEGold, boxrule=0.5pt, arc=2pt,
      left=3pt, right=3pt, top=3pt, bottom=3pt, before skip=0pt, after skip=0pt,
      equal height group=FEnum\the\FEnumgroup, valign=top, halign=center]
    {\color{FEMaroon}\bfseries\Large #1\par}\vspace{2pt}{\scriptsize #2\par}
  \end{tcolorbox}
  \end{minipage}}
\newcommand{\infonumbers}[8]{%
  \global\advance\FEnumgroup by 1
  \infonum{#1}{#2}\hfill\infonum{#3}{#4}\hfill\infonum{#5}{#6}\hfill\infonum{#7}{#8}\par}

% Lab parts -------------------------------------------------------------
\newenvironment{infolab}{\scriptsize\begin{itemize}[leftmargin=1.1em, itemsep=1pt, topsep=1pt, parsep=0pt]}{\end{itemize}}

% Verdict band ----------------------------------------------------------
\newcommand{\infoverdict}[2]{% #1 = where the evidence stands, #2 = what investors can use
  \begin{tcolorbox}[enhanced, colback=white, colframe=FEMaroon, boxrule=0.8pt, arc=3pt,
      leftrule=4pt, left=7pt, right=7pt, top=4pt, bottom=5pt, before skip=0pt, after skip=0pt]
    {\color{FEMaroon}\bfseries\small Where the evidence stands\par}\vspace{2pt}
    {\scriptsize #1\par}\vspace{4pt}
    {\color{FEMaroon}\bfseries\small What investors can use\par}\vspace{2pt}
    {\scriptsize #2\par}
  \end{tcolorbox}}
\fi
\begin{FEinfographic}{ch:20}{do machine-learning models find return predictability that linear models miss, and does it survive costs?}

\infopanel{1}{What machine learning adds}{%
\infocards
  {Regularization}{Shrinkage and the elastic net keep many inputs from fitting noise.}
  {Nonlinearity}{Trees and neural networks capture interactions among familiar signals.}
  {Validation}{Tuning on a held-out period controls overfitting within one design, not across many designs.}}

\infopanel{2}{Where the evidence stands: what published studies found}{%
\begin{infoscore}
\scoreitem{y}{Yes}{Out-of-sample forecasts in broad samples}
\scoreitem{p}{Weakly}{Hidden structure}
\scoreitem{p}{Partly}{Value-weighted portfolios}
\scoreitem{p}{Weakly}{Hard-to-arbitrage stocks excluded, after costs}
\scoreitem{y}{Yes}{A retail-dominated market (China)}
\scoreitem{o}{Disputed}{Very large models for market timing}
\end{infoscore}}

\infopanel{3}{Key numbers}{%
\infonumbers
  {2.45}{equal-weighted Sharpe ratio (Gu, Kelly, Xiu)}
  {1.35}{the same, value-weighted}
  {149}{portfolios in the lab}
  {57\%}{AI fund gain 2017--2022, against 84\% for the S\&P 500}}

\infopanel{4}{Lab: Machine learning on a portfolio panel}{%
\begin{infolab}
\item \textbf{A.} Five models from OLS-3 to a neural network: out-of-sample $R^2$ and long--short portfolios.
\item \textbf{B.} Large portfolios only, before and after 2000, and net of costs.
\item \textbf{C.} Sensitivity to the random seed and input importance. \textbf{D.} The virtue of complexity with random features.
\end{infolab}}

\infosummary
  {In broad samples of U.S. stocks, regularized nonlinear models forecast returns better out of sample than linear models with a few characteristics. But they rely on familiar signals, and their gains are concentrated where trading is costly.}
  {Machine learning is a better blender, not a crystal ball: insist on value-weighted, after-cost results and expect decay after publication.}
\end{FEinfographic}
